\documentclass[11pt,a4paper]{article}
\usepackage[T1]{fontenc}
\usepackage{lmodern,microtype,amsmath,amssymb,amsthm}
\usepackage[margin=25mm,headheight=15pt]{geometry}
\usepackage{booktabs,longtable,array,tabularx,enumitem,xcolor,listings,fancyhdr}
\usepackage[hyphens]{url}
\usepackage[colorlinks=true,linkcolor=blue!45!black,citecolor=blue!45!black,urlcolor=blue!45!black]{hyperref}
\definecolor{ink}{HTML}{17364A}
\definecolor{accent}{HTML}{147D83}
\definecolor{pale}{HTML}{F2F6F8}
\hypersetup{pdftitle={Certified Prime Extensions and Uniform Coefficient Bounds},pdfauthor={Keio},pdfsubject={Exact primality certificates and partial results on Philemon's conjecture}}
\pagestyle{fancy}\fancyhf{}
\fancyhead[L]{\small\color{ink}Certified Prime Extensions}
\fancyhead[R]{\small Keio | 8 November 2026}
\fancyfoot[C]{\thepage}
\setlength{\parskip}{4pt}\setlength{\parindent}{0pt}
\setlength{\emergencystretch}{2em}
\newtheorem{theorem}{Theorem}[section]
\newtheorem{proposition}[theorem]{Proposition}
\newtheorem{corollary}[theorem]{Corollary}
\newtheorem{lemma}[theorem]{Lemma}
\theoremstyle{definition}\newtheorem{definition}[theorem]{Definition}
\theoremstyle{remark}\newtheorem{remark}[theorem]{Remark}
\newcommand{\NN}{\mathbb N}
\newcommand{\ord}{\operatorname{ord}}
\newcommand{\lpf}{P^{-}}
\newcommand{\ph}{\varphi}
\lstset{basicstyle=\ttfamily\scriptsize,breaklines=true,columns=fullflexible,
 keepspaces=true,showstringspaces=false,frame=single,rulecolor=\color{black!15},
 backgroundcolor=\color{pale},numbers=left,numberstyle=\tiny\color{gray},
 numbersep=6pt,xleftmargin=12pt,framexleftmargin=8pt,aboveskip=9pt,belowskip=9pt}
\title{\color{ink}\bfseries Certified Prime Extensions\\and Uniform Coefficient Bounds\\[6pt]
 \large\normalfont Philemon's totient conjecture, exact certificates,\\and a Euclidean prime generator}
\author{Keio}
\date{8 November 2026}
\begin{document}
\maketitle
\begin{abstract}
We identify the giant-prime computations in manuscripts [129] and [130] of
the Verifiable Mathematical Reasoning Corpus, pin their source snapshot,
and independently replay the seven-node certificate chain ending at the
18,886-digit prime $r$. A new search gives
\[
 P=2347\,2^{8192}r+1,
\]
a 21,356-digit prime. Exact modular exponentiation and a greatest common
divisor calculation certify this integer by Pocklington's theorem; the
primality of the old corpus endpoint is an actual dependency of this proof.

We prove Philemon's original totient implication for every coefficient
$1\leq a\leq64$, with no upper bound on $m$, using an explicit weighted
factor-count obstruction and a finite classification of possible coefficient
patterns. An explicit exceptional set has at most $\lfloor X/10\rfloor$
members among the first $X$ coefficients; outside this set, the original
implication holds for every admissible $m$. Thus the proved coefficient
set has lower density at least $9/10$. The unrestricted conjecture is not
settled. We also investigate
Euclidean prime generators. A corollary of the theorems of Booker--Pomerance
and Munsch--Shparlinski proves eventual uniqueness of the new prime factor
in OEIS A359504, while an exact finite computation checks its early terms.
Neither this eventual statement nor the giant-prime computation is a full
resolution of that original conjecture. We make no claim
of a world record, historical first discovery, new general primality algorithm, or
solution satisfying all six requested open-problem success gates.
The complete primality verifier is embedded in this single source file.
\end{abstract}
\vspace{3pt}
\noindent\textbf{Keywords.} Certified prime generation; Pocklington; Proth;
Euclidean prime generators; squarefree smooth integers; reproducible computation.
\newpage

\section{Findings, provenance, and scope}
The starting corpus was read at Git commit
\begin{center}\small\ttfamily c30d40119134a59eda6fa5cc6331f27c8e67d034.\end{center}
The commit timestamp is 6 November 2026, 16:24:52 UTC. The repository's own
README identifies it as a collection of research drafts, not peer-reviewed
or formally verified results. Accordingly, descriptions in the old files
were treated as source claims until the relevant certificate was replayed.
The exact sources are listed in the references \cite{corpus129,corpus130}.

\begin{table}[ht]
\centering\small
\begin{tabularx}{\linewidth}{@{}lXr@{}}
\toprule
Source & Explicitly certified endpoint & Decimal digits\\
\midrule
Corpus [129] & $s=5936p_5+1$ & 18,885\\
Corpus [130] & $r=62988p_5+1$ & 18,886\\
This computation & $P=2347\,2^{8192}r+1$ & 21,356\\

\bottomrule
\end{tabularx}
\caption{Constructed endpoints in the audited experiments. The two old
endpoints $s$ and $r$ are siblings; both have parent $p_5$.
These are not world-record claims.}
\label{tab:records}
\end{table}

The exact definition of $p_5$ is given in Section~\ref{sec:chain}.
The new integer $P$ improves the number of decimal digits in the audited
constructed endpoint by $21,356-18,886=2,470$, approximately $13.08\%$.
This comparison concerns the completed certificate extensions in [129]
and [130]. It excludes the much larger pre-existing Mersenne prime used
as a benchmark input in [130]; $P$ is not larger than that input.
The phrase ``new'' in this report means newly found and certified in the
present computation. It does not assert that no person has previously
considered the same integer.

Corpus [129] also records an interrupted search among
$2ms^2+1$, $1\leq m\leq10^6$, of approximately 37,775 digits. Its durable
log summary contains 2,813 completed Fermat rejections and no certified
prime from that search. The exact tested parameters are not included in
the retrieved source, so their number cannot be used to skip the first
2,813 parameters in a future search. The present search uses a different
family and retains its own machine-readable journals.

The work has two distinct outcomes. Extending the completed corpus
certificate chain has a new unconditional answer. The request to solve an original named
open problem under gates G1--G6 does not have a completed answer here.
Section~\ref{sec:gates} records the precise distinction.

\section{The classical certificate machinery}
\label{sec:certificates}
The following results are included to make the computational proof
self-contained. They are classical; neither changing the coefficient nor
reusing a previously certified factor creates a new primality theorem
\cite{pocklington1914,bls1975,caldwell}.

\begin{theorem}[Pocklington certificate]\label{thm:pocklington}
Let $N>1$, and let $F\mid N-1$ be a completely factored positive integer
with $F^2>N$. For every distinct prime $q\mid F$, suppose that an integer
$a_q$ satisfies
\begin{equation}\label{eq:pocklington}
 a_q^{N-1}\equiv1\pmod N,
 \qquad
 \gcd\bigl(a_q^{(N-1)/q}-1,N\bigr)=1.
\end{equation}
Then $N$ is prime.
\end{theorem}
\begin{proof}
Let $\ell$ be any prime divisor of $N$. The first congruence implies that
$a_q$ is a unit modulo $\ell$ and its order divides $N-1$. The gcd
condition means that this order does not divide $(N-1)/q$. Therefore its
$q$-adic valuation is exactly $v_q(N-1)$; in particular
$q^{v_q(F)}\mid\ell-1$. Applying this reasoning for every prime dividing
$F$ gives $F\mid\ell-1$. Thus $\ell>F>\sqrt N$. A composite integer has
a prime divisor at most its square root, a contradiction.
\end{proof}

\begin{corollary}[One known prime factor]\label{cor:linear}
Let $q$ be a proven prime and $N=cq+1$ with $q^2>N$. If
\[
 a^{N-1}\equiv1\pmod N,\qquad\gcd(a^c-1,N)=1,
\]
then $N$ is prime. No factorization of $c$ is needed.
\end{corollary}

Repeated factors must be counted with their full multiplicity. For
$N=kq^m+1$, one may take $F=q^m$ and use the reduced exponent
$kq^{m-1}$ in \eqref{eq:pocklington}. In particular, for a quadratic lift
the reduced exponent is $kq$, not $k$. Arbitrarily increasing $m$ is
possible at the level of the certificate; it does not guarantee that the
resulting candidate is prime or that the search is inexpensive.

\begin{theorem}[Proth certificate]\label{thm:proth}
Let $K$ be positive and odd, $n\geq1$, and $K<2^n$. If
$N=K2^n+1$ and
\begin{equation}\label{eq:proth}
 a^{(N-1)/2}\equiv-1\pmod N,
\end{equation}
then $N$ is prime.
\end{theorem}
\begin{proof}
Squaring the congruence gives the full Fermat congruence. Since $N$ is
odd, the gcd of the reduced power minus one with $N$ is
$\gcd(-2,N)=1$. Take $F=2^n$ in Theorem~\ref{thm:pocklington}; the strict
bound $K<2^n$ implies $F^2>N$.
\end{proof}

\section{Replay of the existing corpus chain}\label{sec:chain}
Define
\begin{align}
 p_0&=14^{512}-14^{256}+1,\label{eq:seed}\\
 p_i&=k_i p_{i-1}^{\,2}+1\quad(1\leq i\leq5),\label{eq:chain}\\
 (k_1,k_2,k_3,k_4,k_5)&=(438,13696,5494,13258,22330),\notag\\
 r&=62988p_5+1.\label{eq:r}
\end{align}
The certificate starts only from the elementary primes $2$ and $7$,
which the embedded verifier proves by trial division.

\begin{table}[ht]
\centering\small
\begin{tabular}{@{}lrrll@{}}
\toprule
Node & Digits & Bits & Factored part $F$ & Prime factor : base\\
\midrule
$p_0$ & 587 & 1,950 & $2^{256}7^{256}$ & $2:11$, $7:2$\\
$p_1$ & 1,177 & 3,908 & $p_0^2$ & $p_0:2$\\
$p_2$ & 2,357 & 7,829 & $p_1^2$ & $p_1:2$\\
$p_3$ & 4,718 & 15,670 & $p_2^2$ & $p_2:2$\\
$p_4$ & 9,439 & 31,354 & $p_3^2$ & $p_3:2$\\
$p_5$ & 18,882 & 62,722 & $p_4^2$ & $p_4:2$\\
$r$ & 18,886 & 62,738 & $p_5$ & $p_5:2$\\
\bottomrule
\end{tabular}
\caption{All seven rows were independently replayed in this computation.
Every required full-power residue and gcd equals $1$.}
\label{tab:chain}
\end{table}

For the seed, put $x=14^{256}$. Then $p_0-1=x(x-1)$ and
$x^2-p_0=x-1>0$, so the known factors $2^{256}7^{256}$ suffice.
For each quadratic child, the verifier checks the exact bound
$p_{i-1}^4>p_i$. For the final linear child it checks $p_5^2>r$.
Every factor used in a later certificate must already belong to the
verifier's set of proven primes.

The canonical SHA-256 of $r$, computed from its decimal ASCII digits
without a trailing newline, is
\begin{lstlisting}[numbers=none]
ca4c8532ec289a68d0fe9769f87f3797aa560f9f1a141866da17b917c11c1081
\end{lstlisting}
This agrees with the fingerprint in corpus [130]. The hash checks integer
identity; primality comes from the order argument and arithmetic checks.

\section{A 21,356-digit prime using the old certificate}\label{sec:newprime}
\begin{theorem}[Newly certified direct extension]\label{thm:newprime}
With $r$ defined by \eqref{eq:seed}--\eqref{eq:r}, the integer
\begin{equation}\label{eq:newprime}
 \boxed{P=2347\,2^{8192}r+1}
\end{equation}
is prime. It has $21,356$ decimal digits and $70,941$ bits.
\end{theorem}
\begin{proof}
The replay in Section~\ref{sec:chain} proves that $r$ is prime.
The exact computations recorded below and reproduced by the embedded
verifier give
\begin{align}
 r^2&>P,\label{eq:newsize}\\
 2^{P-1}\bmod P&=1,\label{eq:newfermat}\\
 \gcd\bigl(2^{2347\,2^{8192}}-1,P\bigr)&=1.\label{eq:newgcd}
\end{align}
Corollary~\ref{cor:linear} with $q=r$ and $c=2347\,2^{8192}$ proves the
claim. The digit count and bit length are computed from the reconstructed
integer, rather than estimated from floating-point logarithms.
\end{proof}

\subsection{Exact identity and arithmetic replay}
The decimal SHA-256 is
\begin{lstlisting}[numbers=none]
6f4975bbc4efe37c474e15c24dcf3950b130a0338cbb03220550730ad73d4ad8
\end{lstlisting}
The first and last 100 decimal digits are:
\begin{lstlisting}[numbers=none]
First 100:
19009340918408164175883324388920747198517797379045
05369366949958056427946949933422137930810720308457
Last 100:
37652103594107303035768225277297523950490109469258
78577840894258697187649178902755261635260357541889
\end{lstlisting}
The direct replay used gmpy2 2.3.2 with GMP 6.3.0. On this host, while a
separate search was running, the full power in \eqref{eq:newfermat} took
$20.1567$ seconds; the reduced power and gcd in \eqref{eq:newgcd} took
$2.9247$ seconds. These are observed wall times, not complexity bounds or
performance guarantees on another machine.

The complete decimal expansion is included in the companion verification
archive. It is unnecessary for the definition: equations
\eqref{eq:seed}--\eqref{eq:newprime} specify every digit uniquely, and the
single-file verifier reconstructs it without downloading any data.



\section{Search design and retained evidence}\label{sec:search}
The direct continuation considered
\[
 L_k=k\,r\,2^{8192}+1,\qquad
 k\in\{1,3,5,\ldots,1,999,999\}.
\]
There were exactly $10^6$ coefficients before sieving. For each odd prime
$\ell\leq10^8$, let $C\equiv r2^{8192}\pmod\ell$. If $C=0$, that
prime cannot divide any $L_k$. Otherwise the unique prohibited class is
\[
 k\equiv-C^{-1}\pmod\ell.
\]
Writing $k=k_0+2j$ gives the equivalent index condition
\[
 j\equiv(-C^{-1}-k_0)2^{-1}\pmod\ell.
\]
This is an exact sieve of a linear family. It tests $5,761,455$ primes
including $2$ and leaves $60,985$ coefficients. The parity check is
automatic because every candidate is odd.

Three single-thread workers tested survivors using OpenPFGW 4.1.8, with
GWNUM 31.5, in base-two probable-prime mode \cite{pfgw}. The survivor
$k=2347$ has zero-based sorted index $67$ in the retained candidate list.
It passed after $43.3$ seconds of the parallel run. At that moment $67$
tests had completed; after the already-running tests finished, the durable
journals held $69$ completed tests: one pass and $68$ software rejections.
The process ended at $44.1472$ seconds. The completion order of parallel
workers is not the numerical ordering of the coefficients.

The search result was only a screening pass. It became the prime in
Theorem~\ref{thm:newprime} after the separate exact computations
\eqref{eq:newsize}--\eqref{eq:newgcd}. The final proof does not rely on
floating-point FFT roundoff behavior or on the program printing ``PRP''.
The software rejection records describe what that search did; they are
not needed for the primality proof or a claim of a globally first prime.

An earlier, separate Proth search used $K2^{131072}+1$ with $K=202611080000001+2j$, $0\leq j<10^6$. It was stopped after 285 completed software rejections, with no passing candidate. Its journals are retained separately from the next search.

A second Proth search used $K2^{131113}+1$ with $K=202611+2j$, $0\leq j<10^6$. Sieving through $10^8$ left 60,875 coefficients. Seven single-thread workers used the OpenPFGW $N-1$ test. The execution ended without a terminal summary. Reconciliation of the durable journals finds 3,429 completed tests, all software rejections, with no passing or unrecognized result. The last periodic checkpoint recorded 3,425 completions after 1463.1 seconds; later completed records are included in the reconciled count. Unlogged candidates are not classified. This interrupted search supplies no additional prime theorem.

The sieve was checked on independent small-range enumerations of both
families. The downloaded OpenPFGW archive matched the publisher's SHA-256,
\begin{lstlisting}[numbers=none]
aadf885e2d6489866bb7eeea2ba3f75a3fa02e679f1d6a86fd2632e042a99c99
\end{lstlisting}
This provenance check is separate from the mathematical certificate.
The archive accompanying this report preserves the search program, sieve
source, parameter files, journals, and exact verification output. It does
not redistribute the third-party executables.

\section{A bounded audit of the much larger old candidates}
Corpus [130] also discusses
\[
 q=2^{136279841}-1,\qquad
 U=28q^2+1,\qquad V=66q^3+1.
\]
The source gives $82,048,642$ and $123,072,962$ decimal digits, respectively.
These were candidate integers, not certified successes. In particular,
benchmarking a modular square did not evaluate the complete Fermat
congruences that their proposed certificates would require.

The present additional audit checks every prime $\ell\leq10^9$ by
computing $q_\ell=(2^{136279841}-1)\bmod\ell$ and testing
\[
 28q_\ell^2+1\pmod\ell,\qquad66q_\ell^3+1\pmod\ell.
\]
No prime in that range divides either candidate. The run covered
$50,847,534$ primes, including the parity check, and took $8.77406$
seconds on one core. The complete bounded-audit code is reproduced in
Appendix~\ref{app:giant}. Products stay below the $64$-bit unsigned limit
when the enforced bound is at most $10^9$.

This is exactly a finite small-factor exclusion. It does not establish
that $U$ or $V$ is prime, and it does not change the certified maximum in
Table~\ref{tab:records}. No primality assumption about $q$ is needed for
this modular audit.

\section{An original named problem close to the certificate method}
\label{sec:philemon}
The most directly related named target located in the review is
Conjecture 4.1 of Philemon \cite{philemon2021}:
\begin{quote}
Let $n=am+1$, where $a\geq1$, $m>1$, and $p=\lpf(m)$ is the least prime
divisor of $m$. If $a<p$ and $m\mid\ph(n)$, then $n$ is prime.
\end{quote}
This formulation predates the present work and retains the original
quantifiers. The author's version 5, dated 10 April 2021, explicitly
states it as a conjecture. The same paper proves the prime-power case
$m=p^e$ earlier. Original online discussions contain additional partial
arguments, including a case with two distinct prime factors in $m$
\cite{philemonMO,philemonMSE}. No established unrestricted resolution was
located in the sources examined; that search result is not a proof of
historical novelty.

\subsection{Why the giant-prime certificate does not settle it}
Pocklington's gcd condition works separately at every prime divisor
$\ell$ of $n$. It forces a sufficiently large factor of the order modulo
\emph{each} $\ell$, and thereby rules out a small divisor.
The aggregate condition $m\mid\ph(n)$ is different: factors of $m$ may
be supplied by different terms in
\[
 \ph(n)=\prod_{\ell^e\parallel n}\ell^{e-1}(\ell-1).
\]
For example, an order condition that is merely nontrivial modulo $n$
does not automatically make its reduced-power gcd equal to one.
Replacing the original totient condition by the local gcd conditions
would change the problem and would invoke an existing theorem.

An elementary necessary condition illustrates the remaining difficulty.
If a composite counterexample existed, put $t=\ph(n)/m$. Then
\[
 1\leq t<a<p,\qquad\gcd(n,m)=1.
\]
For every prime $u\mid m$, the inequality $t<p\leq u$ implies
$v_u(\ph(n))=v_u(m)$. This constrains the allocation of large prime
factors among the terms $\ell-1$, but does not force all of them into a
single term. Such valuation reductions already occur in the original
discussion; they are not claimed as new results here.

% Standalone section fragment for the root manuscript.
% Required bibliography keys: philemon2021, philemonMSE.
% Required environments: lemma, theorem, corollary, proof.

\section{A factor-count obstruction for Philemon's conjecture}
\label{sec:philemon-obstruction}

Philemon's Conjecture~4.1 asks whether
\begin{equation}\label{eq:philemon-original}
 m>1,\qquad n=am+1,\qquad 1\leq a<p:=\operatorname{lpf}(m),
 \qquad m\mid\varphi(n)
 \quad\Longrightarrow\quad n\text{ is prime}.
\end{equation}
Here $\operatorname{lpf}(m)$ is the least prime divisor of $m$.
The prime-power case and a two-prime case precede the present work
\cite{philemon2021,philemonMSE}.  The factor allocation and product-expansion
strategy used below also have precedents in the 2020 discussion
\cite{philemonMSE}, including the answers by Haran and user140242.
We include the full justification because the distinction between
$\varphi(n)$ and the orders modulo individual prime factors is essential.
The weighted tail estimate and final factor-count bound are derived here,
without a claim of historical priority.  The argument does
not settle the unrestricted implication \eqref{eq:philemon-original}.

\begin{lemma}[Exact allocation in a hypothetical counterexample]
\label{lem:philemon-allocation}
Suppose the hypotheses in \eqref{eq:philemon-original} hold and $n$ is
composite.  Then $m$ is odd.  There are integers $C\geq1$, $s\geq1$,
$c_i\geq2$, and $m_i\geq p$ such that
\begin{equation}\label{eq:philemon-allocation}
 n=C\prod_{i=1}^s(c_i m_i+1),\qquad
 \prod_{i=1}^s m_i=m,\qquad
 U:=C\prod_{i=1}^s c_i<a.
\end{equation}
Every $c_i$ is even.  The numbers $c_i m_i+1$ are precisely the distinct
prime factors of $n$ greater than $p$, each with multiplicity one; every
prime factor of $C$ is less than $p$.  In particular, $C<a$ and
$U\geq C2^s$.  Moreover,
\begin{equation}\label{eq:philemon-tail}
 a-U+\frac1m
 =U\left\{\prod_{i=1}^s\left(1+\frac1{c_i m_i}\right)-1\right\}
 <\frac{Us}{2p-s}.
\end{equation}
\end{lemma}

\begin{proof}
Write $\varphi(n)=bm$.  Since $n$ is composite,
$\varphi(n)<n-1=am$, and hence $1\leq b<a<p$.  Thus
$\gcd(b,m)=1$.  Also $\gcd(n,m)=1$ because $n=am+1$.
If $m$ were even, then $p=2$ and $a=1$, contradicting $b<a$.
Therefore $m$ is odd and $p\geq3$.

Call a prime divisor $\ell$ of $n$ active if
$\gcd(\ell-1,m)>1$.  Any active prime exceeds $p$.  If any prime
$\ell>p$ occurred in $n$ with multiplicity at least two, then
$\ell\mid\varphi(n)$ and $\ell\nmid m$ would give
$\ell\mid b$, contradicting $b<p<\ell$.
If a prime $\ell>p$ were inactive, then
$\ell-1\mid\varphi(n)$ and $\gcd(\ell-1,m)=1$ would give
$\ell-1\mid b$, again impossible.  No prime divisor of $n$ equals
$p$, since $\gcd(n,m)=1$.  Consequently, the active primes are exactly
the prime divisors of $n$ greater than $p$, all with multiplicity one.
List them as $\ell_1,\ldots,\ell_s$ and put
$C=n/(\ell_1\cdots\ell_s)$.  All prime factors of $C$ are less than
$p$, so $\gcd(\varphi(C),m)=1$.  The condition $m>1$ and
$m\mid\varphi(n)$ ensures $s\geq1$.

For each $i$ define the full part of $\ell_i-1$ supported on the prime
divisors of $m$ by
\[
 m_i=\prod_{t\mid m,\ t\ \mathrm{prime}}t^{v_t(\ell_i-1)},
 \qquad c_i=\frac{\ell_i-1}{m_i}.
\]
Because $\gcd(b,m)=\gcd(\varphi(C),m)=1$ and
$\varphi(n)=\varphi(C)\prod_i(\ell_i-1)=bm$, the $t$-adic
valuations on both sides show that $\prod_i m_i=m$ exactly.
Every active $m_i$ is a nontrivial divisor supported on primes at least
$p$, so $m_i\geq p$.  Each $m_i$ is odd, whereas $\ell_i-1$ is even;
therefore $c_i\geq2$ is even.

Expanding $C\prod_i(c_i m_i+1)$ gives
$n=Um+\Delta$, where $\Delta>1$ in the composite case.  Indeed, if
$s\geq2$ the nonleading positive terms already sum to more than one.
If $s=1$, then $\Delta=C$; the alternative $C=1$ would make
$n=\ell_1$ prime.  Comparison with $n=am+1$ now gives
$(a-U)m=\Delta-1>0$, proving $U<a$ and $C<a$.

Dividing the product formula for $n$ by $m$ proves the equality in
\eqref{eq:philemon-tail}.  Since $c_i m_i\geq2p$,
\[
 \prod_i\left(1+\frac1{c_i m_i}\right)
 \leq\left(1+\frac1{2p}\right)^s
 <\frac1{1-s/(2p)}.
\]
The last inequality follows by comparing the binomial expansion with
the geometric series, using $\binom{s}{j}\leq s^j$.
Its denominator is positive: $p>a>U\geq2^s$ implies $s<2p$.
Subtracting one and multiplying by $U$ proves the strict inequality
in \eqref{eq:philemon-tail}.
\end{proof}

\begin{theorem}[Parity-sensitive factor-count obstruction]
\label{thm:philemon-factor-count}
In every composite counterexample to \eqref{eq:philemon-original}, the
notation in Lemma~\ref{lem:philemon-allocation} satisfies the stronger
weighted bound
\begin{equation}\label{eq:philemon-weighted-gap}
 a-U<\sum_{i=1}^s\frac1{c_i}\leq\frac{s}{2}.
\end{equation}
In particular,
\begin{equation}\label{eq:philemon-quantitative-gap}
 2(a-U)<s,
 \qquad\text{equivalently}\qquad s\geq2(a-U)+1.
\end{equation}
Consequently, a composite counterexample must have
at least five distinct prime factors greater than $p$ if $n$ is odd,
and at least three such factors if $n$ is even.
\end{theorem}

\begin{proof}
Use Lemma~\ref{lem:philemon-allocation} and let $h=a-U$ and
$S=\sum_i1/c_i$.  Since $S\leq s/2<a<p$, all the following denominators
are positive.  The elementary symmetric terms in
$\prod_i(1+1/(c_i p))$ are bounded by the corresponding powers of
$S/p$, so comparison with a geometric series gives
\[
 \prod_i\left(1+\frac1{c_i m_i}\right)
 \leq\prod_i\left(1+\frac1{c_i p}\right)
 <\frac1{1-S/p}.
\]
The equality in \eqref{eq:philemon-tail}, followed by $p>a$, yields
\[
 h<\frac{US}{p-S}<\frac{US}{a-S}
   =\frac{(a-h)S}{a-S}.
\]
Multiplication by $a-S$ and cancellation give $ah<aS$, and hence
$h<S\leq s/2$.  This proves both
\eqref{eq:philemon-weighted-gap} and
\eqref{eq:philemon-quantitative-gap}.

The integer $U$ is even.  If $n$ is odd, then $a$ is even because
$m$ is odd; thus $h$ is a positive even integer, so $h\geq2$ and
$s\geq5$.  If $n$ is even, then $a$ and $h$ are odd; thus $h\geq1$
and $s\geq3$.
\end{proof}

\begin{corollary}
Writing $\Omega(m)$ for the number of prime factors of $m$ counted with
multiplicity, a composite counterexample has $\Omega(m)\geq5$ when
$n$ is odd and $\Omega(m)\geq3$ when $n$ is even.  In the notation of
Lemma~\ref{lem:philemon-allocation}, the corresponding lower bounds are
$a\geq32C+2$ and $a\geq8C+1$, respectively.
\end{corollary}

\begin{proof}
Each $m_i>1$ and $\prod_i m_i=m$, so $\Omega(m)\geq s$.
Also $a=U+h\geq C2^s+h$; apply the theorem and the parity bounds on $h$.
\end{proof}

\begin{corollary}[All coefficients up to $32$]
\label{cor:philemon-small-coefficient}
The original implication \eqref{eq:philemon-original} is true for
every integer $1\leq a\leq32$, with no bound on $m$.
\end{corollary}

\begin{proof}
Suppose instead that a composite counterexample has $a\leq32$.
If $n$ is odd, the preceding corollary gives
$a\geq32C+2\geq34$, a contradiction.
If $n$ is even, then $C$ is even and $s\geq3$.  Since
$U=C\prod_i c_i<a\leq32$, the alternatives $s\geq4$, $C\geq4$,
or any $c_i\geq4$ are all impossible: each would give $U\geq32$.
Therefore $C=2$, $s=3$, $c_1=c_2=c_3=2$, and $U=16$.
The strict bound $2(a-U)<s=3$ forces $a-U=1$, so $a=17$.

Now $p>17$ and every prime factor of every $m_i$ is at least $p$;
thus $3\nmid m_i$.  The number $\ell_i=2m_i+1$ is a prime greater
than $p$, so it is not divisible by $3$.  These two facts force
$m_i\equiv2\pmod3$ and $\ell_i\equiv2\pmod3$ for all three
indices.  Consequently the product formula for $n$ gives
$n=2\ell_1\ell_2\ell_3\equiv1\pmod3$, whereas
$n=17m+1=17m_1m_2m_3+1\equiv2\pmod3$.
This contradiction proves the claim.  No primality assumption on
the individual $m_i$ was used.
\end{proof}

% Suggested bibliography entries, to merge into the manuscript bibliography:
% \bibitem{philemon2021}
% A. S. Philemon, Primality of numbers of the form $ap^k+1$,
% arXiv:2005.02327v5 (2021), Conjecture 4.1.
% \url{https://arxiv.org/abs/2005.02327v5}.
%
% \bibitem{philemonMSE}
% Haran, Partial answer to ``Prove that there are no composite integers
% $n=am+1$ such that $m\mid\varphi(n)$'', Mathematics Stack Exchange,
% 7 October 2020, edited 2 November 2020.
% \url{https://math.stackexchange.com/questions/3843281/}.

% Section fragment: append after the factor-count theorem in
% research/prime_strategy/philemon_obstruction.tex.
% This is a uniform coefficient-range result, not a resolution of the
% unrestricted conjecture, and no historical-priority claim is made.

\subsection{A uniform coefficient range through 64}
\label{sec:philemon-64}

\begin{corollary}[All coefficients up to $64$]
\label{cor:philemon-coeff64}
The original implication \eqref{eq:philemon-original} holds for every
integer $1\leq a\leq64$, with no upper bound on $m$.
\end{corollary}

\begin{proof}
Suppose, to the contrary, that $n$ is composite and $a\leq64$.
Use the exact decomposition in Lemma~\ref{lem:philemon-allocation}, and
put $h=a-U$.  Thus $U<a\leq64$, $U\geq C2^s$, all $c_i$ are even,
and \eqref{eq:philemon-quantitative-gap} gives $2h<s$.
The integers $m_i$ are odd, and every prime divisor of every $m_i$ is
at least $p>a$.  Distinct active primes with the same coefficient
$c_i=c_j$ have $m_i\ne m_j$.

\paragraph{Complete coefficient classification.}
If $n$ is odd, then $C$ is odd, $h$ is positive and even, and $s\geq5$.
The inequality $2^s\leq U<64$ forces $s=5$.
Now $U=C\prod_i c_i$ is a positive multiple of $32$ below $64$, so
$U=32$, $C=1$, and all five $c_i$ equal $2$.
The inequality $2h<5$ forces $h=2$, hence $a=34$.

If $n$ is even, then $C$ is even, $h$ is positive and odd, and $s\geq3$.
Here $2^{s+1}\leq U<64$, so $s=3$ or $s=4$.
For $s=4$, $U$ is a positive multiple of $32$ below $64$;
therefore $U=32$, $C=2$, all $c_i=2$, $h=1$, and $a=33$.
For $s=3$, $U$ is a positive multiple of $16$ below $64$, and $2h<3$
gives $h=1$.  Thus $(U,a)$ is one of $(16,17)$, $(32,33)$, and
$(48,49)$.  Factoring
\[
 \frac{U}{16}=\frac{C}{2}\prod_{i=1}^3\frac{c_i}{2}
 \in\{1,2,3\}
\]
gives precisely the following cases, up to permutation of the $c_i$:
\begin{equation}\label{eq:philemon-64-cases}
\begin{array}{c|c|c|c|c|c}
 a&C&s&(c_1,\ldots,c_s)&U&h\\ \hline
 17&2&3&(2,2,2)&16&1\\
 33&2&3&(2,2,4)&32&1\\
 33&4&3&(2,2,2)&32&1\\
 33&2&4&(2,2,2,2)&32&1\\
 34&1&5&(2,2,2,2,2)&32&2\\
 49&2&3&(2,2,6)&48&1\\
 49&6&3&(2,2,2)&48&1
\end{array}
\end{equation}
In particular, this classification includes both the three-factor case
$C=4$ and the four-factor case $C=2$ at $a=33$.

\paragraph{Four exclusions modulo $3$.}
In every row where all $c_i=2$, one has $p>17$, hence $3\nmid m_i$.
Since $2m_i+1$ is prime and exceeds $p$, it is not divisible by $3$.
Consequently $m_i\equiv2\pmod3$ and $2m_i+1\equiv2\pmod3$.
The two formulas for $n$ therefore require
\begin{equation}\label{eq:philemon-64-mod3}
 C2^s\equiv a2^s+1\pmod3.
\end{equation}
For $(a,C,s)=(17,2,3),(33,4,3),(33,2,4),(34,1,5)$, the left and
right residues are, respectively, $(1,2),(2,1),(2,1),(2,0)$.
These four rows are impossible.  The row $(49,6,3)$ survives this
particular congruence and is treated separately below.

\paragraph{Monotone bounds for the remaining rows.}
For positive coordinates $t_i$ define
\[
 R_U(\boldsymbol c;\boldsymbol t)
 =U\left\{\prod_{i=1}^s\left(1+\frac1{c_i t_i}\right)-1\right\}.
\]
This expression decreases in each coordinate $t_i$.
The equality in \eqref{eq:philemon-tail} says
\begin{equation}\label{eq:philemon-64-monotone}
 h+\frac1m=R_U(\boldsymbol c;\boldsymbol m).
\end{equation}
All remaining rows have $h=1$, so any coordinatewise upper bound on
the right side that is less than $1$ gives a contradiction.

For $a=33$, $C=2$, and $\boldsymbol c=(2,2,4)$, one has $p\geq37$.
Order the first two coordinates so that $m_1<m_2$.
The value $m_1=37$ would give the composite active factor $75$.
Every integer from $38$ through $40$ has a prime factor below $37$,
so $m_1\geq41$; oddness and distinctness then give $m_2\geq43$.
Also $m_3\geq37$.  Hence
\begin{equation}\label{eq:philemon-64-bound33}
 R_{32}((2,2,4);\boldsymbol m)
 \leq R_{32}((2,2,4);(41,43,37))
 =\frac{64466}{65231}<1.
\end{equation}

Both remaining rows have $a=49$ and therefore $p\geq53$.
An integer $m_i<53^2$ supported on primes at least $53$ must itself
be prime.  Among integers $d$ in the interval $[53,113]$ for which
both $d$ and $2d+1$ are prime, the possibilities are exactly
\begin{equation}\label{eq:philemon-64-small-values}
 53,\quad83,\quad89,\quad113.
\end{equation}
This small enumeration is elementary: the other primes in the interval
and proper factors of $2d+1$ are listed here:
\[
\begin{array}{c|rrrrrr}
 d&59&61&67&71&73&79\\
 \text{proper factor of }2d+1&7&3&5&11&3&3
\end{array}
\]
\[
\begin{array}{c|rrrrr}
 d&97&101&103&107&109\\
 \text{proper factor of }2d+1&3&7&3&5&3
\end{array}
\]
The four retained numbers $107,167,179,227$ are prime by trial division
through their square roots.  Integers $d\leq113$ that are not prime
cannot be any $m_i$, because their least prime factors are below $53$.

For $C=2$ and $\boldsymbol c=(2,2,6)$, order the two coordinates with
coefficient $2$.  The preceding enumeration gives
$m_1\geq53$, $m_2\geq83$, and $m_3\geq53$.  Consequently
\begin{equation}\label{eq:philemon-64-bound49a}
 R_{48}((2,2,6);\boldsymbol m)
 \leq R_{48}((2,2,6);(53,83,53))
 =\frac{209366}{233147}<1.
\end{equation}

Finally, consider $C=6$ and $\boldsymbol c=(2,2,2)$.
Order all three coordinates increasingly.  If they are not exactly
$(53,83,89)$, their coordinatewise lower bounds are $(53,83,113)$:
the first three eligible values below $113$ are precisely those in
the exceptional triple.  Thus in every nonexceptional case,
\begin{equation}\label{eq:philemon-64-bound49b}
 R_{48}((2,2,2);\boldsymbol m)
 \leq R_{48}((2,2,2);(53,83,113))
 =\frac{477402}{497087}<1.
\end{equation}
The exceptional triple gives, exactly,
\[
 m=53\cdot83\cdot89=391511,
 \qquad n=6\cdot107\cdot167\cdot179=19191306.
\]
But
\begin{equation}\label{eq:philemon-64-exception}
 n-1-49m=19191305-19184039=7266\ne0,
\end{equation}
so it does not satisfy the required identity $n=49m+1$.
Every row of \eqref{eq:philemon-64-cases} is excluded, proving the
corollary.
\end{proof}

The bound in Corollary~\ref{cor:philemon-coeff64} is on the coefficient
$a$, not on the size of $m$ or $n$.  The argument therefore covers
infinitely many possible inputs in each of the $64$ coefficient classes.
It remains a partial result for \eqref{eq:philemon-original}; neither
the remaining coefficients nor historical priority are settled here.

% Append after the factor-count theorem and the a <= 64 corollary.
% Uses labels from philemon_obstruction.tex and philemon_64.tex.
% The density is a density of coefficient classes, uniformly in m.
% No assertion of historical novelty or of a complete conjecture proof.

\subsection{A coefficient set of lower density at least 9/10}
\label{sec:philemon-density}

The quantitative factor-count obstruction also proves the original
implication for an explicit infinite set of coefficients.  Define, for
each positive integer $h$,
\[
 M_h=
 \begin{cases}
  2^{2h+1},&h\text{ even},\\
  2^{2h+2},&h\text{ odd},
 \end{cases}
 \qquad
 \mathcal E_h=\{h+tM_h:t=1,2,3,\ldots\},
 \qquad
 \mathcal E=\bigcup_{h\geq1}\mathcal E_h.
\]
The positive-multiple condition $t\geq1$ is part of the definition.
In particular, the residue representative $a=h$ is not included merely
by belonging to that residue class.

\begin{theorem}[Uniform coefficient-density bound]
\label{thm:philemon-density}
For every coefficient $a\notin\mathcal E$, the original implication
\eqref{eq:philemon-original} holds for every $m>1$ satisfying
$a<\operatorname{lpf}(m)$.  For each positive integer $X$,
\begin{equation}\label{eq:philemon-density-finite}
 \#\bigl(\mathcal E\cap\{1,\ldots,X\}\bigr)
 \leq\left\lfloor\frac{X}{10}\right\rfloor.
\end{equation}
Consequently, at least $X-\lfloor X/10\rfloor$ coefficient classes
among $1,\ldots,X$ satisfy the conjectured implication uniformly in
$m$.  The explicit set $\mathbb Z_{\geq1}\setminus\mathcal E$ has
lower asymptotic density at least $9/10$.
\end{theorem}

\begin{proof}
Suppose a composite counterexample exists, and retain the notation in
Lemma~\ref{lem:philemon-allocation}.  Put $h=a-U>0$.
Theorem~\ref{thm:philemon-factor-count} gives $s\geq2h+1$.
Since every $c_i$ is even, $2^s\mid U$.

If $a$ is even, then $n$ and $C$ are odd, while $h$ is even.
It follows that
\[
 2^{2h+1}\mid U=a-h.
\]
If $a$ is odd, then $n$ and $C$ are even, while $h$ is odd.
The additional factor of two in $C$ gives
\[
 2^{2h+2}\mid U=a-h.
\]
Thus in either case $M_h\mid a-h$.  Because $U=a-h>0$, the quotient
is a positive integer; therefore $a\in\mathcal E_h\subseteq
\mathcal E$.  Contraposition proves the first assertion, with all
quantifiers over $m$ unrestricted.

For a fixed $h$ and $X$, the number of terms in $\mathcal E_h$ up to
$X$ is
\[
 \max\left\{0,\left\lfloor\frac{X-h}{M_h}\right\rfloor\right\}
 \leq\frac{X}{M_h}.
\]
Only finitely many such sets meet $\{1,\ldots,X\}$, and their union
has size at most the sum of their sizes.  Summing the convergent bounds
over all $h$ gives
\begin{align*}
 \#\bigl(\mathcal E\cap\{1,\ldots,X\}\bigr)
 &\leq X\left(
     \sum_{j=1}^{\infty}2^{-(4j+1)}
     +\sum_{j=0}^{\infty}2^{-(4j+4)}\right)\\
 &=X\left(\frac1{30}+\frac1{15}\right)=\frac{X}{10}.
\end{align*}
The left side is an integer, giving
\eqref{eq:philemon-density-finite}.  Taking a lower limit of the
complementary proportions proves the density assertion.
\end{proof}

Membership in $\mathcal E$ is a finite calculation for any specified
$a$: a relevant $h$ must satisfy $h+M_h\leq a$.  The set
$\mathcal E$ contains only coefficients not certified by this density
argument; membership does not assert the existence of a counterexample.
Corollary~\ref{cor:philemon-coeff64} additionally settles every
coefficient up to $64$, including those that lie in $\mathcal E$.
Theorem~\ref{thm:philemon-density} concerns coefficient classes, not
the frequency of primes or a proportion of arbitrary pairs $(a,m)$.
The unrestricted conjecture remains unresolved by these arguments.

\subsection{Restricted unconditional primality certificates}
\label{sec:philemon-restricted-certificate}

The preceding results give an unconditional instance of the certificate
proposed conditionally in Philemon's Theorem~4.2
\cite{philemon2021}.  The order argument itself is the earlier one;
the point is the explicit coefficient range on which its premise has
now been established here.

\begin{corollary}[An unconditional instance of the extended Lucas criterion]
\label{cor:philemon-restricted-certificate}
Let $m>1$, $1\leq a<\operatorname{lpf}(m)$, and $N=am+1$.
Assume either $a\leq64$ or $a\notin\mathcal E$.
Suppose that for every prime $q\mid m$ there is an integer $b_q$
such that
\begin{equation}\label{eq:philemon-restricted-witnesses}
 b_q^{N-1}\equiv1\pmod N,
 \qquad
 b_q^{(N-1)/q}\not\equiv1\pmod N.
\end{equation}
Then $N$ is prime.
\end{corollary}

\begin{proof}
The first congruence makes $b_q$ a unit modulo $N$.
Let $d_q=\operatorname{ord}_N(b_q)$.
Then $d_q\mid N-1$ while $d_q\nmid(N-1)/q$.
Because $a<\operatorname{lpf}(m)$, one has $\gcd(a,m)=1$.
If $q^e\Vert m$, it follows that $q^e\Vert N-1$ and
$q^e\mid d_q\mid\varphi(N)$.  Combining these divisibilities
for the distinct primes dividing $m$ gives $m\mid\varphi(N)$.
Apply Corollary~\ref{cor:philemon-coeff64} or
Theorem~\ref{thm:philemon-density}, as appropriate.
\end{proof}

The noncongruence in \eqref{eq:philemon-restricted-witnesses} is taken
modulo $N$ itself.  This corollary is restricted to its stated
coefficient set; it does not justify omitting a Pocklington gcd check
for arbitrary inputs.  No general speed advantage or historical
priority for a new primality criterion is claimed.


No proof in this report closes that allocation problem in general.
Consequently, Philemon's conjecture cannot be used as an assumption in
the new-prime proofs. The certificates in Sections~\ref{sec:certificates}
and \ref{sec:newprime} are unconditional without it.

\section{A named conjecture on a greedy Euclidean prime sequence}
\label{sec:brennan}

This section records a separate theoretical consequence and an exact bounded
experiment.  The theoretical result is explicitly a corollary of existing
theorems on smooth squarefree representatives.  It does not establish the
all-step assertion in the original conjecture, and it does not depend on the
size of the certified prime constructed elsewhere in this paper.

\subsection{The original sequence and conjecture}

Joel Brennan's OEIS entry A359504, dated January 3, 2023, defines the following
sequence \cite{oeis359504}.  Put $S_0=\varnothing$.  For a finite set $S$ of
primes, write
\[
  \mathcal D(S)=\left\{\prod_{q\in T}q:T\subseteq S\right\},
\]
where the empty product is $1$.  Thus $\mathcal D(S)$ consists of the
squarefree positive integers whose prime factors belong to $S$.  Recursively,
define
\begin{align*}
 P_{n+1}&=\min\{d\in\mathcal D(S_n):
       d+1\text{ has a prime divisor outside }S_n\},\\
 a_{n+1}&=\min\{q:q\text{ is prime},\ q\mid P_{n+1}+1,\ q\notin S_n\},\\
 S_{n+1}&=S_n\cup\{a_{n+1}\}.
\end{align*}
These definitions preserve the original order of the two minimizations:
first minimize the product, and then select its smallest new prime factor.
They do not minimize the new prime across all available products.

The construction continues indefinitely.  Indeed, the product of all primes
in $S_n$ is an eligible product, since every prime factor of that product
plus one lies outside $S_n$.  Each term is a new prime.  The first four terms
are $2,3,7,5$, with selected products $1,2,6,14$ respectively.

\paragraph{Brennan's conjecture.}
At every step, $P_{n+1}+1$ has exactly one distinct prime divisor outside
$S_n$.  The original entry states ``at most one''; eligibility supplies
``at least one,'' so these formulations are equivalent.  Multiplicity of
the new factor is not restricted by the original assertion.

\subsection{Two published inputs}

A positive integer is $p$-smooth when every prime factor is at most $p$.
Booker and Pomerance \cite[Theorem~1]{booker2017} prove that, for every prime
$p\notin\{5,7\}$, every residue class modulo $p$ has a positive squarefree
$p$-smooth representative.  Applied to the residue $-1$, such a representative
has all its prime factors strictly below $p$.

Munsch and Shparlinski \cite[Theorem~1.4]{munsch2020} prove the stronger
asymptotic size bound
\[
 M(p)\le p^{3/2+o(1)},
\]
where $M(p)$ is the smallest bound permitting squarefree $p$-smooth
representatives for all residue classes.  In particular, for each fixed
$\varepsilon>0$, every sufficiently large prime $p$ admits such a
representative $d$ of $-1$ with $d\le p^{3/2+\varepsilon}$.

Booker and Pomerance's Corollary~7 concerns a different selection rule: it
minimizes the new prime over all products.  The argument below transfers
their representative theorem to Brennan's product-first rule.

\begin{corollary}[Every prime occurs]
The sequence A359504 contains every prime exactly once.

\end{corollary}\begin{proof}
Distinctness follows from the definition.  Suppose that a prime never
occurs, and let $p$ be the least such prime.  The displayed initial terms
give $p\ge11$, handling both exceptional moduli in the published theorem.
All primes smaller than $p$ occur, so there is an index $N$ after which
all of them belong to $S_n$.

Choose a squarefree integer $d$ with all prime factors below $p$ and
$d\equiv-1\pmod p$, using Booker--Pomerance.  For every $n\ge N$, the
integer $d$ belongs to $\mathcal D(S_n)$ and remains eligible, because
$p\mid d+1$ and $p\notin S_n$.  Therefore $P_{n+1}\le d$ and
$a_{n+1}\le d+1$ for every $n\ge N$.  This would give infinitely many
distinct primes bounded by the fixed integer $d+1$, a contradiction.
\end{proof}

Let $\mu_n$ denote the least prime outside $S_n$.  The first corollary implies
$\mu_n\to\infty$: each fixed finite set of primes is eventually contained
in $S_n$.

\begin{corollary}[Eventual uniqueness, with multiplicity one]
For each $\varepsilon>0$, there is an index $N_\varepsilon$ such that
\[
 P_{n+1}\le \mu_n^{3/2+\varepsilon},\qquad
 a_{n+1}\le \mu_n^{3/2+\varepsilon}+1
 \quad(n\ge N_\varepsilon).
\]
Consequently, for all sufficiently large $n$, the integer $P_{n+1}+1$
has exactly one prime divisor outside $S_n$, and that prime occurs to
exponent one.

\end{corollary}\begin{proof}
For sufficiently large $n$, Munsch--Shparlinski provides a squarefree
$\mu_n$-smooth integer $d_n\equiv-1\pmod{\mu_n}$ with
$d_n\le \mu_n^{3/2+\varepsilon}$.  Since $\mu_n\nmid d_n$, all prime
factors of $d_n$ are below $\mu_n$ and hence belong to $S_n$.
Thus $d_n$ is eligible, and the two bounds follow from the selection rule.

Fix $0<\varepsilon<1/2$.  As $\mu_n\to\infty$, these bounds give
$P_{n+1}+1<\mu_n^2$ eventually.  Every prime outside $S_n$ is at least
$\mu_n$.  If $P_{n+1}+1$ had two distinct new prime divisors, or if a new
prime divisor occurred at least twice, it would be at least $\mu_n^2$.
Both possibilities are excluded.  Eligibility guarantees the existence
of a new prime divisor, completing the proof.\end{proof}

The proof also gives an exact local criterion: at any step, an eligible
squarefree integer $d\equiv-1\pmod{\mu_n}$ with prime factors below $\mu_n$
and $d+1<\mu_n^2$ certifies uniqueness and exponent one at that step.
This criterion is sufficient; its failure is not a counterexample.

\subsection{Exact finite experiment and its scope}

We implemented a smallest-prime-factor sieve and an availability queue for
all squarefree products $d\le B$.  For each such product, a counter stores
how many of its distinct prime factors have not yet appeared.  Adding a
new prime decrements these counters on the corresponding squarefree
multiples.  Products whose counters reach zero enter a min-heap.  A
popped product is factored after adding one, and the distinct factors
outside the current prime set are counted.

This handles decreases in the selected products, such as the transition
from $14$ to $10$ after the first four terms.  At every successful step,
the least eligible product in the bounded queue is also the least
eligible product globally, because every omitted product exceeds $B$.
Already inspected products cannot acquire a new absent factor when the
known-prime set grows.  These invariants justify the resulting sequence
prefix; exhaustion of the bounded queue is only a stopping condition for
the experiment.

With $B=100{,}000{,}000$, the run generated the first $3{,}795{,}778$
sequence terms and inspected $58{,}574{,}992$ available squarefree subset products.
No selected product plus one had two distinct new prime factors.  The
largest generated prime was $99{,}999{,}971$.  The recorded running time
was $22.75$ seconds in the execution environment used for this paper;
this is an observation rather than a hardware-independent performance
claim.  The first $5{,}000$ terms were compared against the complete
official OEIS b-file, with no mismatches; the common $5{,}000$th term is
$77{,}291$.  The finite search checks distinct new factors, which is the
property in Brennan's conjecture.

\subsection{The remaining logical gap}

The eventual-uniqueness corollary proves that any failures of Brennan's conjecture occur among
finitely many steps.  The proof does not supply a numerical index after
which failures are impossible.  The finite experiment cannot cover an
unspecified exceptional range.  Combining the asymptotic theorem with
the displayed search therefore does not prove the original assertion
for every step.

For the requested success conditions, the target retains its published
definition, but the established statement has the weaker quantifier
``all sufficiently late steps.''  Thus it does not satisfy the required
full-scope G2/G4 closure.  The theoretical argument is a corollary of
published representation bounds, so it supplies no independent claim
of the methodological novelty demanded by G5.  No priority or G6 claim
is made.  The explicit outcomes are the two corollaries and the bounded
verification, with the original all-step conjecture left unresolved.


\section{Success-gate audit and remaining conclusions}\label{sec:gates}
\begingroup\small\setlength{\tabcolsep}{4pt}
\begin{longtable}{@{}>{\raggedright\arraybackslash}p{0.07\linewidth}>{\raggedright\arraybackslash}p{0.33\linewidth}>{\raggedright\arraybackslash}p{0.54\linewidth}@{}}
\toprule
Gate & Required success criterion & Result of this investigation\\
\midrule\endhead
G1 & A previously named problem with a prior original statement &
Philemon Conjecture 4.1 and OEIS A359504 are concrete prior targets. A new
large prime by itself is not such a solved conjecture.\\
G2 & Preserve the original proposition and quantifiers &
The targets are stated unchanged. The coefficient restrictions and
eventual uniqueness are explicit partial results; neither replaces an
original unrestricted quantifier.\\
G3 & The target was actually unresolved &
The primary sources examined still state the targets conjecturally. No
claim is made that a search can rule out all unpublished or overlooked
prior work.\\
G4 & An unconditional complete proof of the original target &
Complete proofs are supplied for the specific prime certificates, $a\leq64$, and the explicit
coefficient set of lower density at least $9/10$, all uniformly in $m$.
Not achieved for either unrestricted original
conjecture. The unspecified eventual cutoff in A359504 is not closed.\\
G5 & A substantive new argument, not a re-expression &
The prime tests are classical. The weighted factor-count estimate works
for a partial range of Philemon's conjecture; A359504 uses existing
analytic theorems. No new method resolving an entire target is claimed.\\
G6 & No prior identical result or essentially identical proof &
No historical-first claim is justified for the integers, algorithms, or
analytic corollaries in this report.\\
\bottomrule
\end{longtable}
\endgroup

The request to extend the completed corpus chain is fully answered.
The partial theorems do not
meet G4 for either original conjecture, and the historical priority required
by G6 has not been established.

\section*{Reproducibility and attribution}
This report was prepared with assistant-supported source review,
mathematical analysis, programming, and exact computation. Classical
theorems are attributed to their sources. The original corpus is licensed
CC BY 4.0; its seed definitions and certificate parameters are reused with
attribution to manuscripts [129] and [130]. The verifier below is
self-contained: the complete source needs no external bibliography,
integer file, image, or custom style to compile. A normal TeX installation
with the packages in the preamble suffices.

\appendix
\section{Complete exact-integer primality verifier}\label{app:verifier}
The program reconstructs the complete dependency chain, proves the small
base primes by trial division, and checks every certificate explicitly.
It uses gmpy2 when available; otherwise Python's exact integers supply the
same mathematical operations with a slower runtime. The checks are
ordinary conditional statements and remain active under Python's
optimization flag. No probable-prime oracle is used.
% EMBEDDED_FILE_BEGIN verify_certificates.py
\begin{lstlisting}[language=Python]
"""Exact primality proofs. Python 3; optional gmpy2 accelerates the same checks.
No input files, prime tables, PRP oracle, or network are used.
"""
import hashlib
import json
import math
import sys
import time

if hasattr(sys, "set_int_max_str_digits"):
    sys.set_int_max_str_digits(0)
try:
    import gmpy2
    Z, modpow, gcd = gmpy2.mpz, gmpy2.powmod, gmpy2.gcd
    backend = "gmpy2 " + gmpy2.version()
except ImportError:
    Z, modpow, gcd = int, pow, math.gcd
    backend = "Python standard-library exact integers"

def need(condition, explanation):
    if not condition:
        raise ArithmeticError(explanation)

known = set()
for q in (2, 7):
    need(all(q % d for d in range(2, math.isqrt(q) + 1)),
         "base-prime trial division failed")
    known.add(Z(q))

def fingerprint(label, n, expected_digits, expected_bits,
                expected_hash=None):
    decimal = str(n)
    digest = hashlib.sha256(decimal.encode("ascii")).hexdigest()
    need(len(decimal) == expected_digits, label + ": digit count")
    need(n.bit_length() == expected_bits, label + ": bit length")
    if expected_hash is not None:
        need(digest == expected_hash, label + ": decimal identity")
    print(json.dumps({"label": label, "status": "CERTIFIED",
          "digits": len(decimal), "bits": n.bit_length(),
          "sha256_decimal": digest}), flush=True)
    return n

def pocklington(label, n, factors):
    need(n > 2 and n % 2, label + ": odd positive integer")
    F = Z(1)
    distinct = set()
    for q, exponent, a in factors:
        need(q in known, label + ": unknown factor primality")
        need(exponent >= 1 and q not in distinct,
             label + ": invalid factor list")
        distinct.add(q)
        F *= q**exponent
    need((n-1) % F == 0, label + ": F does not divide n-1")
    need(F*F > n, label + ": insufficient factored part")
    for q, exponent, a in factors:
        need(1 < a < n, label + ": witness range")
        need(modpow(Z(a), n-1, n) == 1,
             label + ": Fermat congruence failed")
        t = modpow(Z(a), (n-1)//q, n)
        need(gcd(t-1, n) == 1, label + ": gcd failed")
    known.add(n)
    return n

started = time.perf_counter()
x = Z(14)**256
p = pocklington("p0", x*x-x+1,
                [(Z(2), 256, 11), (Z(7), 256, 2)])
fingerprint("p0", p, 587, 1950)
sizes = [(1177,3908), (2357,7829), (4718,15670),
         (9439,31354), (18882,62722)]
for i, (k, size) in enumerate(zip(
        (438,13696,5494,13258,22330), sizes), 1):
    parent = p
    p = pocklington("p"+str(i), k*parent*parent+1,
                    [(parent,2,2)])
    fingerprint("p"+str(i), p, *size)
r = pocklington("r", 62988*p+1, [(p,1,2)])
fingerprint("r", r, 18886, 62738,
    "ca4c8532ec289a68d0fe9769f87f3797"
    "aa560f9f1a141866da17b917c11c1081")
P = pocklington("P", 2347*(Z(2)**8192)*r+1, [(r,1,2)])
fingerprint("P", P, 21356, 70941,
    "6f4975bbc4efe37c474e15c24dcf3950b"
    "130a0338cbb03220550730ad73d4ad8")



print("ALL CERTIFICATES PASSED;", backend,
      "elapsed seconds", round(time.perf_counter()-started, 3))

\end{lstlisting}
% EMBEDDED_FILE_END verify_certificates.py

\section{Source for the bounded giant-candidate audit}\label{app:giant}
Compile with \texttt{g++ -O3 -std=c++17 audit.cpp -o audit} and run
\texttt{./audit 1000000000}. At $2$, both integers are odd; all remaining
primes are generated by an Eratosthenes sieve. Modular exponentiation
uses only exact integer operations.
% EMBEDDED_FILE_BEGIN audit_giant_small_factors.cpp
\begin{lstlisting}[language=C++]
#include <cstdint>
#include <cstdlib>
#include <iostream>
#include <vector>
#include <chrono>

static std::uint64_t power_mod(std::uint64_t a, std::uint64_t e, std::uint64_t m) {
    std::uint64_t r=1;
    while(e) { if(e&1) r=r*a%m; e>>=1; if(e) a=a*a%m; }
    return r;
}

int main(int argc,char**argv) {
    const std::uint64_t bound=argc>1 ? std::strtoull(argv[1],nullptr,10) : 100000000;
    if(bound>1000000000ULL) return 2;
    auto begin=std::chrono::steady_clock::now();
    std::vector<bool> composite(bound/2+1,false);
    for(std::uint64_t p=3;p*p<=bound;p+=2) if(!composite[p/2])
        for(std::uint64_t n=p*p;n<=bound;n+=2*p) composite[n/2]=true;
    std::uint64_t count=1,hits2=0,hits3=0;
    for(std::uint64_t p=3;p<=bound;p+=2) if(!composite[p/2]) {
        ++count;
        std::uint64_t a=power_mod(2,136279841,p);
        std::uint64_t q=(a+p-1)%p;
        std::uint64_t sq=q*q%p;
        if((28*sq+1)%p==0) { ++hits2; std::cout << "N2_FACTOR " << p << " q_mod " << q << std::endl; }
        if((66*(sq*q%p)+1)%p==0) { ++hits3; std::cout << "N3_FACTOR " << p << " q_mod " << q << std::endl; }
    }
    double dt=std::chrono::duration<double>(std::chrono::steady_clock::now()-begin).count();
    std::cout << "bound="<<bound<<" primes_checked="<<count<<" N2_hits="<<hits2<<" N3_hits="<<hits3<<" seconds="<<dt<<std::endl;
}

\end{lstlisting}
% EMBEDDED_FILE_END audit_giant_small_factors.cpp

\section{Source for the exact A359504 finite search}\label{app:frontier}
This program maintains the original priority rule, including products
smaller than the previously selected product that become eligible after
a new prime is inserted. Its invariant and the precise scope of the
finite result are explained in the main text. The plus-sign run is the
one used for the A359504 claim. The optional minus-sign mode is retained
for reproducibility of the exploratory code; no general theorem about
that mode is asserted here.
% EMBEDDED_FILE_BEGIN frontier.cpp
\begin{lstlisting}[language=C++]
#include <algorithm>
#include <chrono>
#include <cstdint>
#include <fstream>
#include <functional>
#include <iostream>
#include <queue>
#include <string>
#include <vector>
using namespace std;
int main(int argc,char**argv){
 const int M=argc>1?stoi(argv[1]):1000000;
 const int sign=argc>2?stoi(argv[2]):1;
 auto start=chrono::steady_clock::now();
 vector<int> spf(M+2),ps;
 for(int i=2;i<=M+1;i++){
  if(!spf[i]){spf[i]=i;ps.push_back(i);}
  for(int p:ps){if(p>spf[i]||1LL*p*i>M+1)break;spf[p*i]=p;}
 }
 vector<int8_t> missing(M+1,0);
 for(int i=2;i<=M;i++){
  int p=spf[i],v=i/p;
  missing[i]=(v%p==0||missing[v]<0)?-1:missing[v]+1;
 }
 vector<bool> seen(M+2,false);
 priority_queue<int,vector<int>,greater<int>> heap;
 heap.push(1);
 int n=0,largest=0,maxP=0,examined=0; bool failure=false;
 ofstream prefix(string("research/named_problem/prefix_")+(sign==1?"plus":"minus")+".csv");
 prefix<<"n,P,new_prime\n";
 auto add=[&](int q,int P){
  n++; largest=max(largest,q);seen[q]=true;
  if(n<=10000)prefix<<n<<","<<P<<","<<q<<"\n";
  if(n<=59)cout<<q<<((n==59)?"\n":" ");
  for(int x=q;x<=M;x+=q){if(missing[x]>=0){missing[x]--;if(missing[x]==0)heap.push(x);}}
 };
 if(sign==-1){add(2,0);add(3,0);}
 while(!heap.empty()){
  int P=heap.top();heap.pop();examined++;maxP=max(maxP,P);
  int x=P+sign;
  vector<int> unseen;
  while(x>1){int p=spf[x]; if(!seen[p])unseen.push_back(p); while(x%p==0)x/=p;}
  if(unseen.size()>1){cout<<"COUNTEREXAMPLE step="<<n+1<<" P="<<P<<" factors=";for(int p:unseen)cout<<p<<",";cout<<"\n";failure=true;break;}
  if(unseen.size()==1)add(unseen[0],P);
 }
 double elapsed=chrono::duration<double>(chrono::steady_clock::now()-start).count();
 cout<<"RESULT M="<<M<<" sign="<<sign<<" terms="<<n<<" largest_prime="<<largest<<" max_P="<<maxP<<" products="<<examined<<" failure="<<failure<<" elapsed_seconds="<<elapsed<<"\n";
 return 0;
}

\end{lstlisting}
% EMBEDDED_FILE_END frontier.cpp

\clearpage
\begin{thebibliography}{99}\small\raggedright
\bibitem{corpus129} Keio,
\emph{Certified Prime Lifting from Known Primes}, corpus manuscript [129],
file \texttt{Certified\_Prime\_Lifting.tex}, snapshot
\texttt{c30d40119134a59eda6fa5cc6331f27c8e67d034}, 2026.
\url{https://github.com/QEEDDD/Verifiable-Mathematical-Reasoning-Corpus}.
\bibitem{corpus130} Keio,
\emph{Certified Prime Lifting and Sparse Modular Squaring}, corpus
manuscript [130], same snapshot, 2026. Repository as in [1].
\bibitem{pocklington1914} H. C. Pocklington,
The determination of the prime or composite nature of large numbers by
Fermat's theorem, \emph{Proceedings of the Cambridge Philosophical Society}
\textbf{18} (1914--1916), 29--30.
\bibitem{bls1975} J. Brillhart, D. H. Lehmer, and J. L. Selfridge,
New primality criteria and factorizations of $2^m\pm1$,
\emph{Mathematics of Computation} \textbf{29} (1975), 620--647.
\bibitem{caldwell} C. Caldwell, PrimePages,
\emph{Primality Proving 3.1: $n-1$ tests and Pepin's tests for Fermats}.
\url{https://t5k.org/prove/prove3_1.html}. Accessed 8 November 2026.
\bibitem{pfgw} OpenPFGW project, Linux release 4.1.8 and command-line
documentation. Accessed 8 November 2026.
\url{https://sourceforge.net/projects/openpfgw/files/pfgw_linux_4.1.8.7z/}.
\url{https://github.com/primesearch/OpenPFGW/blob/master/distribution/pfgwdoc.txt}.
\bibitem{philemon2021} A. S. Philemon,
\emph{Primality of numbers of the form $ap^k+1$},
arXiv:2005.02327v5, 10 April 2021, Conjecture 4.1.
\url{https://arxiv.org/abs/2005.02327}.
\bibitem{philemonMO} A. S. Philemon and contributors,
\emph{Prove that there are no composite integers $n=am+1$ such that
$m\mid\phi(n)$}, MathOverflow question 373497, 2020.
\url{https://mathoverflow.net/questions/373497}.
\bibitem{philemonMSE} A. S. Philemon and contributors,
same question, Mathematics Stack Exchange question 3843281, 2020.
\url{https://math.stackexchange.com/questions/3843281}.
\bibitem{oeis359504} J. Brennan, OEIS A359504, original entry
3 January 2023, with extension by K. Ryde, 10 January 2023.
\url{https://oeis.org/A359504}. Accessed 8 November 2026.
\bibitem{booker2017} A. R. Booker and C. Pomerance,
\emph{Squarefree smooth numbers and Euclidean prime generators}, 2017,
Theorem 1 and Section 5.
\url{https://math.dartmouth.edu/~carlp/nothing8.5.pdf}.
\bibitem{munsch2020} M. Munsch and I. E. Shparlinski,
On smooth square-free numbers in arithmetic progressions,
\emph{Journal of the London Mathematical Society} (2020),
doi:10.1112/jlms.12297; arXiv:1710.04705v3, Theorem 1.4.
\url{https://arxiv.org/abs/1710.04705}.
\end{thebibliography}
\end{document}
